\CGSetHeaderLeft{SERRA DO CUÓ}
\CGTableOfContents
\clearpage

\CGZoneStats{Global}
\pagebreak
\clearpage

%\CGSetHeaderRight{APRESENTAÇÃO}

%\CGTextSingleColumn{Apresentação}{
%Inspirado em grandes encontros da comunidade e no constante desenvolvimento de novos picos de rocha no Brasil, este guia nasce do esforço conjunto para consolidar a escalada na Serra do Cuó, em Campo Grande-RN.

%Mais do que um catálogo, este material tem como objetivo auxiliar a comunidade no reconhecimento do local, organizar o acesso aos diversos setores e, acima de tudo, promover a integração sustentável entre os escaladores e a população local. Todo o desenvolvimento esportivo na região contou, ao longo do tempo, com o fundamental apoio da Prefeitura Municipal de Campo Grande e da Associação de Escaladores do Rio Grande do Norte (AERN).

%Preparamos este guia para que ele seja a base definitiva das suas aventuras na Serra do Cuó. Que todos tenham excelentes escaladas!}

%\CGDoubleTextSection{Primeira Impressão: O Despertar para o Cuó}{
%\textit{Por Flávia dos Anjos}

%Minha primeira impressão da Serra do Cuó foi incrível. A vista de longe é de duas enormes paredes laranja que só me lembram as imagens que vi das tão aclamadas paredes da Catalunha espanhola. Já passavam das 9h30 da manhã e o sol não estava de brincadeira. A trilha era claramente longa e exposta, o que me deixou com um sentimento misto de “eu preciso ir nessa linda parede AGORA!” intercalado com “Eu não quero subir essa trilha de jeito NENHUM!”.

%Para minha sorte (ou não), os amigos Janine Falcão e Allysson Laurentino tinham programado algo bem mais amigável para o pouco tempo que tínhamos disponível. Acontece que, ao longo da inclinada vegetação que se forma aos pés dos dois paredões, a superfície é recoberta de enormes blocos de pedra. Muitos blocos, com todas as faces verticais e negativas, oferecendo inúmeras vias e permitindo que os escaladores escolhessem o lado da sombra e se ajustassem a ela ao longo do dia.

%O que mais me chamou a atenção foi a existência de fendas, muitas fendas. Escalei algumas vias ali e no fim me deparei com uma linha que me chamou muita atenção. Perguntei que via era e o Allysson me confirmou que era um projeto; o conquistador tinha montado a linha, mas não tinha conseguido a tão desejada “cadena”.
%}{}{
%Me interessei e resolvi tentar. O grau da via concentra-se nos dois movimentos iniciais, um \textit{boulder} de calcanhar seguido de um movimento dinâmico. A cadena só saiu depois de um descanso, um lanchinho e um terceiro e último “pega”. Virando o \textit{boulder} a via fica bem amigável, a ser protegida com peças móveis médias.

%Já que o Allysson me garantiu que eu estava levando a “primeira ascensão” (FA), ele me deu o direito de batizar a via, e eu sugeri “Ainda não mijei hoje”.

%Depois de um dia no Cuó eu olhei para trás de novo e voltei a encarar os dois paredões. Refleti sobre o local, sobre aquele projeto sem cadena e terminei com minha segunda e última impressão do Cuó: muito potencial! Uma rocha magnífica, cheia de fendas, com oportunidades para todos os níveis de escalada ou de disposição.}

%\clearpage
%\CGSetHeaderRight{VISÃO GERAL E HISTÓRIA}

%\CGDoubleTextSection{Visão Geral e História local}{
%A palavra "Cuó" é de origem indígena e significa acauã, ou seja, “pássaro devorador de cobras”. É o mesmo termo que deu origem ao nome do município de Caicó/RN, que significa \textit{Queicuó} e, traduzido, quer dizer “rio acauã”. No entanto, segundo o renomado historiador Luís da Câmara Cascudo (1968), Cuó significa divisão ou separação na língua nativa. A Serra do Cuó é separada por um boqueirão, dando-lhe um aspecto típico dividido, em uma região historicamente povoada pelos cariris.

%Hugolino D’Oliveira, recenseador no ano de 1940, já se referia à Serra do Cuó detalhando sua geografia: um pico com cerca de 428 metros de altura na ponta oriental da serra, formando o boqueirão que o separa do restante do maciço.
%}{}{
%Hoje a Serra do Cuó não abriga mais seus antigos habitantes. Uma nova geração se levantou e transformou a montanha em um formidável palco para os esportes de aventura. Suas trilhas, blocos e enormes paredões de pedra acolhem agora pessoas de vários lugares do Brasil e do mundo em busca de superação, liberdade e contato direto com a natureza.}

%\CGTextSingleColumn{A História da Escalada na Serra do Cuó}{
%As escaladas na Serra do Cuó tiveram seu início oficial em outubro de 2013, com a conquista das primeiras vias no Talhado dos Americanos. Brito Filho e Patrício Fernandes deram a primeira investida com a abertura de linhas esportivas. A via \textit{Boca Torta}, com proteção fixa, foi a pioneira, seguida imediatamente por outras duas fendas tradicionais com proteção em móvel e parada fixa.

%Após 2015, devido a alguns problemas de acesso ao setor dos blocões originais e às trilhas que partiam de lá, as atividades de desenvolvimento na serra sofreram um período de esfriamento. Mais recentemente, no entanto, o ímpeto de conquistas retornou com muita força, unindo novas gerações a antigos escaladores. Atualmente, o grande foco do montanhismo local está nos setores mais novos: Paraíso e Catedral.}
